% Paper 4, §13 — Scaling interpretability and frontier-2026 snapshot
% Thesis: each of the five methods has a different scaling envelope. SAEs
% scale on documented L(C), L(N,K) laws; probes are graceful by construction;
% attribution graphs scale to production but expensively and only inside the
% originating lab; manual circuit hunts and full activation patching do not
% yet scale to frontier-2026 architectures (MLA / fine-grained MoE / NSA).
% This section is the compute-and-coverage map of the paper.

\section{Scaling interpretability and frontier-2026 snapshot}
\label{sec:scaling-snapshot}

The five methods of \S\S\ref{sec:logit-lens-family}--\ref{sec:circuits}
were originally exhibited on small open models: GPT-2 small (124M
parameters) for \citet{wang2022-ioi}'s indirect-object-identification
circuit, GPT-2 XL (1.5B) for \citet{meng2022-rome}'s ROME tracing,
Pythia 410M--12B for \citet{belrose2023-tuned-lens}'s tuned-lens
benchmarks. The 2024--2026 question --- and the operational question for
anyone who wants interpretability evidence on a model people actually
deploy --- is which of these methods carry to 2--27B open-weights models
(the Gemma 2 / LLaMA 3 / Qwen 2.5 family) and to production frontier
models (Claude 3.5 Haiku, GPT-4 class). The answer is method-dependent,
not uniform: SAEs scale on documented laws, probes scale by construction,
attribution graphs scale to production but only inside Anthropic, and
manual circuit hunts scale not at all. This section is the
compute-and-coverage map; \S\ref{sec:scaling-snapshot-table} consolidates
it into a single snapshot table that becomes the figure-of-record for
the paper.

\subsection{What scales}
\label{sec:scaling-yes}

\paragraph{SAEs scale on documented laws.}
\citet{gao2024-topk-saes} report clean scaling laws for sparse
autoencoders trained with a TopK activation: a compute--loss law
$L(C)$ at fixed sparsity, and a width--sparsity law $L(N, K)$ at fixed
compute, where $N = M$ is the dictionary size and $K$ the per-token
$L_0$ \citep[\S 1, \S 3]{gao2024-topk-saes}\footnote{KB excerpt:
\texttt{kb/excerpts/gao2024-topk-saes.md\#sec-1-headline}.}. The headline
demonstration is a 16M-latent SAE trained on GPT-4 residual-stream
activations with $\sim 7\%$ dead latents under the AuxK
mitigation~\citep[\S 2.4]{gao2024-topk-saes}\footnote{KB excerpt:
\texttt{kb/excerpts/gao2024-topk-saes.md\#sec-2-4}.}; the law is
predictive enough to guide budget allocation between $N$ and training
tokens at the design stage rather than empirically. \citet{gemma-scope-2024}
exercise the architectural cousin --- JumpReLU SAEs of
\citet{rajamanoharan2024-jumprelu} --- at the open-weights end of the
same scaling envelope: more than $400$ SAEs spanning every layer and
sublayer of Gemma~2 2B and 9B (residual stream, MLP output, attention
output) plus selected layers of Gemma~2 27B, with widths
$\{16.4\text{K}, 65.5\text{K}, 131\text{K}, 1\text{M}\}$, totaling
$>\!30$M learned features and roughly $20\%$ of GPT-3's pretraining
compute~\citep[\S 1]{gemma-scope-2024}\footnote{KB excerpt:
\texttt{kb/excerpts/gemma-scope-2024.md\#sec-1}.}. The CC-BY-4.0 release
on HuggingFace plus the Neuronpedia interactive demo
\deepencite{Neuronpedia integration: \texttt{neuronpedia.org/gemma-scope},
not in papers.json} converted SAE-based interpretability from an
Anthropic-internal capability into an externally reproducible one.

\paragraph{Probes scale gracefully.}
A linear probe is trained on cached activations (so the model itself is
called once to produce activations, then the probe runs on a CPU);
mass-mean probing (\S\ref{sec:mass-mean}) is closed-form
$\mathbf{d}_{\mathrm{MM}} = \boldsymbol{\mu}_+ - \boldsymbol{\mu}_-$ and
requires no training at all~\citep{marks-tegmark-2023-truth}. Probe
\emph{cost} is dominated by the activation cache, not the probe fit;
the activation cache scales linearly with dataset size and is
embarrassingly parallel across prompts. The same property carries the
direction-ablation operator
$\mathbf{h} \mapsto \mathbf{h} - (\mathbf{h}^\top \hat{\mathbf{d}}_y)\,\hat{\mathbf{d}}_y$
and the addition operator $\mathbf{h} \mapsto \mathbf{h} + \alpha\,
\hat{\mathbf{d}}_y$ of \S\ref{sec:ablate-add}: both are $O(D)$ per token
and survive whatever architectural innovations the underlying model
adds, because they act on the residual stream rather than on any
particular sublayer.

\paragraph{Attribution graphs scale to production --- expensively, and
inside one lab.}
\citet{lindsey2025-circuit-tracing}'s cross-layer-transcoder framework
reports the only published end-to-end mechanistic-interpretability
result on a production model (Claude 3.5 Haiku). The compute argument
is two-part: per-task cost on the replaced model is one forward pass
plus one backward pass --- cheap relative to activation patching's
$O(N_{\mathrm{components}})$ forward passes per (clean, corrupted) pair
(\S\ref{sec:patching-protocol}) --- but the amortized cost is the
cross-layer transcoder training, estimated at the same order as SAE
training in \citet{gemma-scope-2024} ($\sim\! 20\%$ of the underlying
model's pretraining compute). The Anthropic 2025-06 open release of
the circuit-tracing tooling and its Neuronpedia integration
\deepencite{Anthropic open-source release of circuit-tracer 2025-06,
not in papers.json} converted the methodology from a strictly internal
capability into one external groups can in principle exercise on open
models. As of 2026, no externally-trained attribution-graph result on a
frontier model has been published; the contradiction this raises about
reproducibility is the load-bearing one of \S\ref{sec:contradictions}.

\subsection{What does not scale (yet)}
\label{sec:scaling-no}

\paragraph{Manual circuit hunts.}
The IOI reverse-engineering of \citet{wang2022-ioi} took graduate-student
months of work to identify $26$ heads in $7$ functional classes inside a
$1.6\%$-of-(head, position)-pairs subgraph on a $124$M-parameter model.
The same effort would need to scale roughly with the number of heads
and layers --- $L \cdot H$ on a frontier-2026 model is $\sim 80$--$120$
times that of GPT-2 small --- and the exhaustive case-by-case behavioral
testing that gave the IOI claim its load-bearing status does not have an
obvious automated analog. \citet{conmy2023-acdc} automate the
\emph{discovery} step (iterative path-patching with a KL-divergence
threshold), but the automation produces larger and less interpretable
subgraphs, and the human-in-the-loop labeling that decided ``this head
class is a Name Mover'' on IOI does not transfer.

\paragraph{Full activation patching.}
The three-run protocol of \S\ref{sec:patching-protocol} requires
$O(N_{\mathrm{components}})$ forward passes per (clean, corrupted) pair
to fill a causal-trace heatmap. On GPT-2 XL, $N_{\mathrm{components}}$
is order-$10^4$ at head granularity and $10^5$ at single-(head,
position); on a frontier-2026 model with $\sim 80$ layers, $\sim 64$
heads per layer, and $\sim 8\text{K}$ context, full-coverage patching
is on the order of $10^7$ forward passes per task. Attribution patching
\citep[forum signal, \deepencite{nanda2023-attribution-patching
blog post; not in papers.json}] linearizes the indirect effect via the
first-order Taylor approximation
$\mathrm{IE}_C \approx \langle \nabla_{\mathbf{h}} m,
\Delta \mathbf{h} \rangle$
(\S\ref{sec:patching-attribution}) and reduces this to a small constant
number of passes, but the linearization fails in non-linear regions of
the loss landscape, so it is a discovery aid rather than a faithful
substitute for direct intervention.

\paragraph{Lens coverage of frontier-2026 architectures.}
\citet{belrose2023-tuned-lens}'s tuned-lens benchmarks cover Pythia
410M--12B and GPT-2 / GPT-Neo / BLOOM; \texttt{logitlens4llms-2025}
extends the framework to Qwen-2.5 and LLaMA-3.1
(\S\ref{sec:lens-trajectory}). Tuned lenses for the rest of the
frontier-2026 lineup --- LLaMA 3/4, full Qwen 2.5/3, DeepSeek-V3,
Gemma 3 --- are not all published as of 2026. The relevant
architectural shift for the lens framework is not raw scale (the
distillation loss
$\mathcal{L}(\ell) = \E_{\mathbf{x}}[D_{\mathrm{KL}}(\mathcal{M}_{>\ell}
(\mathbf{h}_\ell) \| \mathrm{TunedLens}_\ell(\mathbf{h}_\ell))]$ is
indifferent to model size) but the substrate: DeepSeek-V3 combines MLA
\citep{deepseek-v2}'s low-rank-compressed KV with fine-grained MoE
\citep{deepseekmoe2024}; NSA \citep{yuan2025} bakes attention sparsity
into pretraining. The residual stream remains well-defined in all of
these, but the assumption that a single per-layer affine translator
$(A_\ell, \mathbf{b}_\ell)$ suffices to absorb basis drift may need
re-examination on a per-architecture basis. The empirical work that
would close this gap --- training tuned lenses across the frontier-2026
lineup and reporting the per-architecture distillation-loss curves ---
is not yet in the literature.

\begin{contradiction}
\textbf{The interpretability-coverage gap is widening, not closing.}
The frontier-2026 architecture snapshot (Paper~1, §13) catalogs MLA,
fine-grained MoE, NSA-sparse attention, sliding-window layers, and
hybrid attention/SSM stacks as the substrate of the deployed lineup.
The lens benchmarks of \citet{belrose2023-tuned-lens} predate this
substrate; the SAEs of \citet{gemma-scope-2024} and the attribution
graphs of \citet{lindsey2025-circuit-tracing} cover specific models in
the lineup but not the lineup as a class; manual circuits and full
activation patching are exhibited only on $\le$1.5B-parameter
GPT-family models. The gap between ``what frontier models do'' and
``what we have interpretability evidence about frontier models doing''
is not narrowing on net: it is widening, because architectural
innovation outpaces the per-architecture interpretability work needed
to catch up. The 2025-06 Anthropic open release of circuit-tracing
tooling and the Gemma Scope CC-BY-4.0 release are corrective forces;
whether they correct fast enough is a load-bearing empirical question
about the field's coverage trajectory.
\end{contradiction}

\subsection{Production interpretability pipelines, 2026}
\label{sec:scaling-pipelines}

Three pipelines are in active use as of 2026, each anchored on one of
the methods above.

\paragraph{Anthropic --- circuit tracer + Neuronpedia.}
The cross-layer-transcoder + Jacobian-attribution pipeline of
\citet{lindsey2025-circuit-tracing}, open-sourced in 2025-06 and
integrated with Neuronpedia for interactive feature browsing
\deepencite{Anthropic 2025-06 open-source release; Neuronpedia
integration; not in papers.json}. The substantive findings on Claude
3.5 Haiku --- forward-planning rhyme features, multi-step factual
chains (``Texas'' $\to$ ``Austin''), refusal-circuit features,
partial-cross-lingual feature overlap --- are reported in
\citet{lindsey2025-biology} and remain Anthropic-internal as evidence
\deepencite{lindsey2025-biology HTML-only; verbatim transcription
pending Phase 2}.

\paragraph{OpenAI --- TopK SAEs and the $L(C), L(N, K)$ scaling
program.}
The $L(C)$ and $L(N, K)$ laws of \citet{gao2024-topk-saes} are the
ongoing methodological backbone, plus an evaluation-metric suite (
$L_0$, reconstruction MSE, downstream-loss recovery, contrastive probe
recovery)~\citep[\S 1 contributions]{gao2024-topk-saes}\footnote{KB
excerpt: \texttt{kb/excerpts/gao2024-topk-saes.md\#sec-1-headline}.}.
The 16M-latent GPT-4 SAE is the existence proof for SAE scaling at
production-frontier scale; the per-feature interpretability work
downstream of that release is not all public.

\paragraph{Google DeepMind --- Gemma Scope + JumpReLU + tuned-lens
benchmarks.}
The JumpReLU SAE family of \citet{rajamanoharan2024-jumprelu}, the
Gemma Scope open release of \citet{gemma-scope-2024}, and the Gemma 2
2B/9B/27B coverage make this the most externally-reproducible pipeline
of the three. The tuned-lens benchmarks of
\citet{belrose2023-tuned-lens}, retargeted across the Gemma family,
provide the cross-method readout layer the lens-as-decoder framing
needs.

\subsection{Snapshot table}
\label{sec:scaling-snapshot-table}

Table~\ref{tab:method-scaling} consolidates the per-method coverage as
of 2026-05. ``Open release'' is the date the method's
canonical-scale-target artifact (SAE weights, lens parameters,
attribution-graph tooling) became externally usable; ``2026
reproduction'' is whether at least one externally-validated result on
a frontier-2026 model exists.

\begin{table}[t]
\centering
\caption{Interpretability methods --- scale, target, and coverage as of
2026-05. \emph{Target} is the largest model the method has been
reported on. \emph{Open release} is the date the canonical
scale-target artifact became externally available; ``\,---\,'' if no
such release. \emph{2026 reproduction} marks whether the
canonical-scale finding has been independently reproduced outside the
originating lab on a frontier-2026 model.}
\label{tab:method-scaling}
\small
\begin{tabularx}{\textwidth}{@{}l l l l l X@{}}
\toprule
Method & Canonical paper & Target model & Scale & Open release &
2026 reproduction \\
\midrule
Logit lens & \citet{nostalgebraist2020-logit-lens} \deepencite{nostalgebraist2020-logit-lens
LessWrong post; not in papers.json} & GPT-2 family & $\le$1.5B &
no fixed artifact (decoding rule) & yes (GPT-2, BLOOM, GPT-Neo);
fails on BLOOM/GPT-Neo per \citet{belrose2023-tuned-lens} \\
Tuned lens & \citet{belrose2023-tuned-lens} & Pythia 12B; GPT-Neo X &
$\le$20B & training script + benchmarks (2023) &
partial: \texttt{logitlens4llms-2025} on Qwen-2.5 / LLaMA-3.1; no
Gemma 3 / DeepSeek-V3 / LLaMA-4 published \\
Mass-mean probe & \citet{marks-tegmark-2023-truth} &
LLaMA-2-13B & 13B & code on GitHub (2023) & yes for
truth, sycophancy variants on LLaMA family; OOD generalization
unresolved (\S\ref{sec:truth-frontier}) \\
ReLU+L1 SAE & \citet{templeton2024-scaling-monosemanticity}
\deepencite{templeton2024-scaling-monosemanticity HTML-only at
transformer-circuits.pub} & Claude 3 Sonnet & frontier &
no public weights (Anthropic-internal) & no \\
TopK SAE & \citet{gao2024-topk-saes} & GPT-4 & 16M-latent & training
recipe + evaluation suite (2024) & yes at $\le$Gemma 2 27B scale;
no GPT-4-class external reproduction \\
JumpReLU SAE & \citet{rajamanoharan2024-jumprelu};
\citet{gemma-scope-2024} & Gemma 2 27B & 1M-latent & Gemma Scope
weights, CC-BY-4.0, HF (2024-08) & yes (most reproducible pipeline) \\
Activation patching & \citet{meng2022-rome};
\citet{zhang2023-apatching} & GPT-2 XL; LLaMA-7B
\deepencite{zhang2023-apatching scale extension; PDF-grounded but no
frontier replication} & 1.5B--7B & toolkits (TransformerLens, etc.)
& yes for $\le$7B; full-coverage at frontier scale infeasible
without attribution patching \\
Manual circuit hunt & \citet{wang2022-ioi};
\citet{conmy2023-acdc} & GPT-2 small & 124M & ACDC code
(2023) & yes on GPT-2 small; not yet on frontier-2026 models \\
Cross-layer transcoder + attribution graph
& \citet{lindsey2025-circuit-tracing} & Claude 3.5 Haiku & frontier &
circuit-tracer tooling (2025-06)
\deepencite{Anthropic open-source release 2025-06; Neuronpedia
integration; not in papers.json} & no external reproduction as of
2026 \\
\bottomrule
\end{tabularx}
\end{table}

The table makes the asymmetry explicit. The methods with the most
favorable scaling --- probes and SAEs --- have the most external
reproduction; the method with the strongest substantive findings on a
production frontier model --- attribution graphs --- has none. Three
of the nine rows depend on tooling or claims that are HTML-only or
released-but-not-yet-reproduced; the remaining six form a defensible
shared methodological substrate.

\begin{intuition}
The five-method picture maps cleanly to the scaling table. Lenses give
you a cheap readout that may not generalize across architectures (no
fixed artifact, partial coverage). Probes give you cheap
correlational evidence with closed-form variants and broad reproduction
(\S\ref{sec:mass-mean}). SAEs give you a feature dictionary at scale
under documented laws (the most reproducible pipeline). Activation
patching gives you causal claims at small scale and degrades to
linearized approximations beyond. Circuit tracing gives you the
strongest claims at frontier scale but only inside the originating
lab. The methodological lesson of \S\ref{sec:cross-method} ---
method-pluralism is principled because methods answer different
questions --- has a scaling-table corollary: \emph{method-pluralism
is also operationally necessary, because no single method covers every
(model, behavior, evidence-class) cell of the matrix}. Returning to
canonical form: the per-method costs are $O(D)$ for probes and lenses,
$O(M D \cdot \text{tokens})$ amortized for SAE training,
$O(N_{\mathrm{components}})$ per-task for activation patching, and
$O(\text{transcoder training} + 1\,\text{fwd} + 1\,\text{bwd})$ for
attribution graphs.
\end{intuition}

\subsection{Forward link}
\label{sec:scaling-forward}

The substrate this section reads against is the frontier-2026
architecture snapshot of \emph{Paper~1, §13}: the MLA / fine-grained
MoE / NSA-sparse-attention / sliding-window / hybrid-SSM stack that
defines what production models look like as of 2026. The
interpretability methods catalogued here have caught up to that
substrate unevenly, and the two most pressing gaps --- per-architecture
tuned lenses for the frontier lineup, and external reproduction of
attribution-graph findings outside Anthropic --- are open empirical
questions whose closure is the work of the next two years. The
contradictions surfaced in this section --- circuit-tracing
reproducibility, lens cross-architecture coverage, the widening gap
between deployed-model architecture and per-architecture
interpretability evidence --- are taken up alongside the
SAE-canonicality and truth-direction contradictions in
\S\ref{sec:contradictions}, where each is paired with the experiment
that would resolve it. The thesis of the paper holds against this
scaling picture: ``\emph{partially read}'' is exactly the right
phrase, and the qualifier is load-bearing across method, scale, and
architecture.
